% ===================== page 1 =====================
\printedpage{}
% head: (aucun ; page de titre non numérotée)

\begin{center}
{\Large\textsf{THE DERIVATORS}}

\bigskip

\textbf{ALEXANDRE GROTHENDIECK}

\bigskip\bigskip

\textbf{Chapter VIII}

\medskip

\textbf{Relative homotopy 1-types: their integration \dots}
\end{center}

\clearpage

% ===================== page 2 =====================
\printedpage{}
% head: (aucun ; page non numérotée)

This text was deciphered and transcribed in \LaTeXe{} by M. Künzer. It was edited by M. Künzer, J. Malgoire, and G. Maltsiniotis. The transcription is as faithful as possible to the manuscript. For the few obvious corrections, or rare comments by the editors, as well as for the original numbering of the manuscript pages, typewriter characters \texttt{[typewriter]} in brackets are used. A question mark in brackets means that one is not sure of the preceding word.

This edition is provisional. Remarks, comments and corrections are welcome. Send a message to:

\textbf{[e-mail address removed]}

\hfill G. Maltsiniotis\hspace*{2cm}

\clearpage

% ===================== page 3 =====================
\printedpage{1}
% head: A. Grothendieck : Dérivateurs, Ch. VIII, Édition des Universités de Montpellier II et Paris VII
% foot: Version 2 février 2002

\msp{page 1}

\section*{1\quad Integration of relative homotopy $1$-types over a category.}

Relative homotopy types over a (small) category $S$ are defined by \textit{stacks} (in categories) on $S^{\wedge}$, which correspond essentially to contrafunctors $S^o \longrightarrow (\mathrm{Cat})$, or again to categories $X$ over $S$ which are Cat-fibered over $S$. (In $(\mathrm{Cat}/S)_{\mathrm{fib}}$, one inverts the arrows which are Hot-equivalences on the fibers.) On the other hand, relative homotopy 1-types over $S$ should be described by \textit{stacks in groupoids} on $S$. If one has a relative homotopy type, defined by
\[ X \longrightarrow S, \qquad\qquad \text{Cat-fibered category over } S, \]
or by
\[
\begin{aligned}
S^o \quad &\longrightarrow \quad \mathrm{Cat}\\
s \quad &\longmapsto \quad X_s\ ,
\end{aligned}
\]
the associated homotopy 1-type is defined by
\[ s \longmapsto \Pi(X_s), \qquad\qquad S^o \longrightarrow (\mathrm{Grp}), \qquad\qquad \text{category of groupoids,} \]
where $\Pi(X_s)$ is the fundamental groupoid of $X_s$. We denote the corresponding fibered category over $S$ by \fbox{$\Pi(X/S)$}. It is true that if
\[ X \ \longrightarrow\ X' \]
is a homomorphism of fibered categories over $S$ which induces a Hot-equivalence on the fibers, it induces a fibered equivalence
\[ \Pi(X/S) \ \longrightarrow\ \Pi(X'/S), \]
but not an isomorphism. There would thus be work to do to define a functor
\[
\underbrace{\mathrm{Hot}(S)}_{\text{\footnotesize\begin{tabular}{c}category of\\ relative homotopy\\ types over $S$\end{tabular}}}
\quad\xrightarrow{\ \xi\,\longmapsto\,\Pi(\xi/S)\ }\quad
\underbrace{\Grp(S)}_{\text{\footnotesize\begin{tabular}{c}category of\\ groupoids\\ fibered over $S$\\ (or categories\\ of stacks in\\ groupoids\\ on $S^{\wedge}$)\end{tabular}}}
\quad .
\]

But this is not my purpose for the moment. Rather, to give a transitivity formula in noncommutative homology, analogous to $R(gf)_* \simeq Rg_*Rf_*$ in cohomology.

\clearpage

% ===================== page 4 =====================
\printedpage{2}
% head: A. Grothendieck : Dérivateurs, Ch. VIII, Édition des Universités de Montpellier II et Paris VII
% foot: Version 2 février 2002

\msp{page 2}

If
\[ f : S \longrightarrow T \]
is a functor, we are going to define a functor
\[ f_! :\ \Grp(S) \ \longrightarrow\ \Grp(T) \]
(direct image of relative homotopy 1-types). But let us first note that if
\[ p :\ X \ \longrightarrow\ S \]
is any $S$-category \textit{whatsoever} (not necessarily Cat-fibered), one nevertheless associates to it a relative homotopy type, via the fibered category $\tilde{X}^{(S)}$, or $\tilde{X}$, having as fibers the $s\backslash X$ (noting that $s' \longrightarrow s$ in $S$ does indeed define $s\backslash X \longrightarrow s'\backslash X$) (\footnote{Having $\tilde{X}$, one defines $\Pi(X/S) \overset{\text{def}}{=} \Pi(\tilde{X}/S)$.}). We have a canonical $S$-functor
\[ X \xrightarrow{\ i\ } \tilde{X}, \qquad\qquad x \longmapsto (x, p(x), \underbrace{p(x)}_{=\,\text{‘}\,s\,\text{’}} \xrightarrow{\ \mathrm{id}\ } p(x)). \]
This is an $S$-Hot-equivalence if $X$ was already fibered. In fact, on each fiber it is a homotopism, for one then defines $\tilde{X}_{s_0} \xrightarrow{\ j\ } X_{s_0}$ in an obvious way, in such a way that
\[
X_{s_0} \underset{j}{\overset{i}{\rightleftarrows}} \tilde{X}_{s_0}, \qquad ji \simeq \mathrm{id}, \qquad ij \underbrace{\xrightarrow{\text{simple homotopy}}}_{\text{\footnotesize\begin{tabular}{l}NB this arrow\\ is not \texttt{[an]}\\ isomorphism\\ in general.\end{tabular}}} \mathrm{id}.
\]
Let us then return to $f : S \longrightarrow T$, and consider a fibered groupoid over $S$,
\[ G \ \longrightarrow\ S, \]
this is not a fibered category over $T$, but one can associate to it
\[ t \longmapsto t\backslash G, \qquad\qquad T^o \longrightarrow (\mathrm{Cat}), \]
hence a fibered category over $T$, of which one takes the relative $\Pi$. But it is simpler to take directly
\[ t \ \longmapsto\ \Pi(t\backslash G), \]
which defines a category fibered in groupoids over $T$.

\clearpage

% ===================== page 5 =====================
\printedpage{3}
% head: A. Grothendieck : Dérivateurs, Ch. VIII, Édition des Universités de Montpellier II et Paris VII
% foot: Version 2 février 2002

\msp{page 3}

Thus we do indeed find a functor (in fact, even a 2-functor of 2-categories)
\[ f_! :\ \Grp(S) \ \longrightarrow\ \Grp(T). \]
If
\[ u :\ G \ \longrightarrow\ G' \]
is a homomorphism of fibered groupoids over $S$, which induces Cat-equivalences on the fibers, then
\[ f_!(u) :\ f_!(G) \ \longrightarrow\ f_!(G') \]
is also a Cat-equivalence on the fibers. (This is formal, from the fact that $f_!$ is a 2-functor, taking into account that if one has an equivalence on the fibers, then there is a quasi-inverse -- this is a general story about fibered categories \dots)

Now let
\[ p :\ X \ \longrightarrow\ S, \]
and let us take the case $G = \Pi(X/S)$. I say that we then have a $T$-equivalence
\[ \boxed{\ f_!(\Pi(X/S)) \ \simeq\ \Pi(X/T),\ } \tag*{$(*)$} \]
in other words, to know $\Pi(X/T)$ (up to $T$-equivalence), it suffices to know $\Pi(X/S)$, and to apply to it the functor $f_!$ of integration of relative homotopy 1-types, for $f : S \longrightarrow T$.

Let us note that we trivially have
\[ \Pi(X/S) \ =\ p_!(e_X), \]
where $e_X$ denotes the trivial fibered $X$-groupoid (all fibers one-point). Thus the preceding formula is written as a transitivity formula
\[ f_!(p_!(\Gamma)) \ \simeq\ (fp)_!(\Gamma), \]
\msp{page 4}
where $\Gamma$ is the trivial fibered groupoid over $X$. This suggests generalizing the formula $(*)$ to
\[ \boxed{\ (fp)_!(\Gamma) \ \simeq\ f_!(p_!(\Gamma))\ } \]
for \textit{every} fibered groupoid over $X$, in a situation
\[ X \xrightarrow{\ p\ } S \xrightarrow{\ f\ } T. \]

This will be formal, if we establish a 2-adjunction property for the functor $f_!$: